\subsection{Solving a System -- Substitution Method} 
Rather than drawing a graph, we can solve a system algebraically. We have two methods for this. The first is the \emph{substitution method}.
\begin{example}
Use substitution to solve the system
\[
\begin{system}
y=3x+5\\
3x+2y=1
\end{system}
\]
\begin{lpic}{}{10}
\end{lpic}
The solution is \blankpair
\end{example}
\begin{enumerate}
\item Solve one equation for one variable. (We'll usually choose substitution when this is \makeblank.)
\item Substitute that into the second equation.
\item You now have an equation in \makeblank[.5in] variable. Solve it.
\item You now have the \makeblank[1in] of one variable.
\item Plug into one of the equations to find the value of the \makeblank[1in] variable.
\item Check by plugging into the other equation.
\end{enumerate}